% Part:sets-functions-relations
% Chapter: size-of-sets
% Section: non-enumerability

\documentclass[../../../include/open-logic-section]{subfiles}

\begin{document}

\olfileid{sfr}{siz}{nen}

\olsection{Himpunan kang \printtoken{S}{nonenumerable}}

\begin{editorial}
  Bagean iki mbuktekake yen $\Bin^\omega$ lan
  $\Pow{\PosInt}$ ora bisa didhaptar, nganggo definisi ing \olref[enm]{sec}.
  Bagean iki dirancang supaya rada luwih dhasar lan luwih rinci
  tinimbang versi ing \olref[enm-alt]{sec}
\end{editorial}

Sawetara himpunan, kayata himpunan $\PosInt$ saka wilangan wutuh positif,
iku tanpa wates. Nganti saiki kita wis weruh tuladha himpunan tanpa wates
kang kabeh !!{enumerable}. Nanging, ana uga himpunan tanpa wates
kang ora duwe sipat iki. Himpunan kaya mangkono diarani \emph{!!{nonenumerable}}.

Kapisan, anane himpunan kang !!{nonenumerable} bisa uga wis nggumunake.
Kanggo saben himpunan~$A$ kang !!{enumerable}, ana
fungsi !!a{surjective} $f \colon \PosInt \to A$. Yen sawijining himpunan
!!{nonenumerable}, ora ana fungsi kaya mangkono. Tegese, ora ana fungsi
kang metakake !!{element}s saka~$\PosInt$ kang cacahe tanpa wates menyang~$A$
lan bisa nyakup kabeh~$A$. Dadi ana !!{element}s saka~$A$ kang
"luwih akeh" tinimbang wilangan wutuh positif kang cacahe tanpa wates.

Kepriye carane mbuktekake yen sawijining himpunan !!{nonenumerable}?
Kudu dituduhake yen ora bisa ana fungsi surjektif kaya mangkono.
Iki padha karo nuduhake yen anggota~$A$ ora bisa dienumerasi ing dhaptar
tanpa wates kang lumaku menyang siji arah. Cara kang paling becik yaiku
nuduhake yen saben dhaptar !!{element}s saka~$A$ mesthi ora ngemot paling
ora siji anggota, utawa ora ana fungsi $f\colon \PosInt \to A$ kang
!!{surjective}. Iki bisa ditindakake nganggo \emph{metodhe diagonal} Cantor.
Saka dhaptar !!{element}s saka~$A$, upamane $x_1$, $x_2$, \dots,
kita mbangun anggota~$A$ liyane kang, amarga cara anggone dibangun,
ora mungkin ana ing dhaptar mau.

Tuladha kapisan yaiku himpunan~$\Bin^\omega$ saka kabeh rerangken
tanpa wates kang ora ana sela lan dumadi saka $0$ lan $1$.

\begin{thm}
\ollabel{thm:nonenum-bin-omega}
$\Bin^\omega$~iku !!{nonenumerable}.
\end{thm}

\begin{proof}
Kanggo mbuktekake kanthi kontradiksi, upamakna $\Bin^\omega$ iku
!!{enumerable}, yaiku ana dhaptar $s_{1}$, $s_{2}$,
$s_{3}$, $s_{4}$, \dots{} saka kabeh !!{element}s ing~$\Bin^\omega$.
Saben $s_i$ dhewe iku rerangken tanpa wates saka $0$ lan~$1$.
Anggota kaping $j$ saka rerangken kaping $i$ ing dhaptar iki kita sebut
$s_i(j)$. Dadi rerangken kaping $i$, yaiku~$s_i$, iku
\[
s_i(1), s_i(2), s_i(3), \dots
\]

Dhaptar iki, bebarengan karo anggota saben rerangken $s_i$,
bisa ditata dadi larikan:
\[
\begin{array}{c|c|c|c|c|c}
& 1 & 2 & 3 & 4 & \dots \\\hline
1 & \mathbf{s_{1}(1)} & s_{1}(2) & s_{1}(3) & s_1(4) & \dots \\\hline
2 & s_{2}(1)& \mathbf{s_{2}(2)} & s_2(3) & s_2(4) & \dots \\\hline
3 & s_{3}(1)& s_{3}(2) & \mathbf{s_3(3)} & s_3(4) & \dots \\\hline
4 & s_{4}(1)& s_{4}(2) & s_4(3) & \mathbf{s_4(4)} & \dots \\\hline
\vdots & \vdots & \vdots & \vdots & \vdots & \mathbf{\ddots}
\end{array}
\]
Label ing sisih kiwa nuduhake nomer rerangken ing dhaptar
$s_1$, $s_2$, \dots; nomer ing sisih ndhuwur menehi label marang
!!{element}s saben rerangken. Upamane, $s_{1}(1)$ iku jeneng
kanggo wilangan, yaiku $0$ utawa~$1$, kang dadi !!{element} kapisan
ing rerangken $s_{1}$, lan sateruse.

Saiki kita mbangun rerangken tanpa wates $\overline{s}$ saka $0$ lan
$1$ kang ora mungkin ana ing dhaptar iki. Definisi
$\overline{s}$ bakal gumantung marang dhaptar $s_1$, $s_2$, \dots.
Saben dhaptar tanpa wates saka rerangken tanpa wates $0$ lan $1$
ngasilake rerangken tanpa wates~$\overline{s}$ kang mesthi ora katon
ing dhaptar mau.

Kanggo netepake $\overline{s}$, kita nemtokake kabeh !!{element}sne,
yaiku nemtokake $\overline{s}(n)$ kanggo kabeh $n \in \PosInt$.
Carane yaiku maca diagonal larikan ndhuwur saka ndhuwur mudhun
(mula diarani "metodhe diagonal"), banjur ngganti saben $1$
dadi $0$ lan saben $0$ dadi~$1$. Kanthi luwih abstrak,
$\overline{s}(n)$ ditetepake dadi $0$ utawa $1$ miturut apa
!!{element} kaping $n$ ing diagonal, yaiku $s_n(n)$, iku $1$ utawa $0$.
\[
\overline{s}(n) =
\begin{cases}
1 & \text{yen $s_{n}(n) = 0$}\\
0 & \text{yen $s_{n}(n) = 1$}.
\end{cases}
\]
Yen luwih seneng rumus tinimbang definisi kanthi kasus,
kita uga bisa netepake $\overline{s}(n) = 1 - s_n(n)$.

Cetha yen $\overline{s}$ iku rerangken tanpa wates saka $0$ lan
$1$, amarga rerangken iki mung mbalikke saben $0$ lan $1$
kang katon ing diagonal larikan. Dadi $\overline{s}$
iku !!a{element} saka~$\Bin^\omega$. Nanging rerangken iki ora bisa
ana ing dhaptar $s_1$, $s_2$, \dots{} Yagene?

Rerangken iki ora bisa padha karo rerangken kapisan $s_1$,
amarga beda karo $s_1$ ing !!{element} kapisan. Apa wae $s_1(1)$,
kita netepake $\overline{s}(1)$ dadi walikane. Rerangken iki uga
ora bisa padha karo rerangken kapindho, amarga $\overline{s}$ beda karo
$s_2$ ing anggota kapindho: yen $s_2(2)$ iku $0$,
$\overline{s}(2)$ iku $1$, lan kosok baline. Mangkono sateruse.

Kanthi luwih pas: yen $\overline{s}$ ana ing dhaptar, mesthi ana
$k$ saengga $\overline{s} = s_{k}$. Rong rerangken padha
yen lan mung yen padha ing saben posisi, yaiku kanggo saben~$n$,
$\overline{s}(n) = s_{k}(n)$. Mligine, kanthi njupuk $n = k$,
kudu kelakon $\overline{s}(k) = s_{k}(k)$. $s_k(k)$ iku
$0$ utawa~$1$. Yen iku $0$, mula $\overline{s}(k)$ kudu~$1$,
miturut cara kita netepake $\overline{s}$. Nanging yen $s_k(k) = 1$,
maneh miturut cara kita netepake $\overline{s}$, $\overline{s}(k) = 0$.
Ing loro-lorone kasus, $\overline{s}(k) \neq s_{k}(k)$.

Kita miwiti kanthi nganggep ana dhaptar !!{element}s saka
$\Bin^\omega$, yaiku $s_1$, $s_2$, \dots{} Saka dhaptar iki kita
mbangun rerangken~$\overline{s}$ kang kabukten ora bisa ana ing dhaptar.
Nanging rerangken iki mesthi dumadi saka $0$ lan $1$ yen kabeh $s_i$
iku rerangken saka $0$ lan $1$, yaiku $\overline{s} \in
\Bin^\omega$. Iki mligine nuduhake yen ora bisa ana dhaptar
\emph{kabeh} !!{element}s saka~$\Bin^\omega$, amarga saka saben
dhaptar kaya mangkono kita bisa mbangun rerangken~$\overline{s}$
kang mesthi ora ana ing dhaptar. Mula anggepan yen ana dhaptar
kabeh rerangken ing~$\Bin^\omega$ nuwuhake kontradiksi.
\end{proof}

\begin{explain}
Cara pambuktian iki diarani "diagonalisasi" amarga nggunakake diagonal
larikan kanggo netepake~$\overline{s}$. Diagonalisasi ora kudu
nggunakake larikan: kanthi gagasan kang padha, kita bisa nuduhake
yen himpunan ora !!{enumerable} sanajan ora ana larikan
lan ora ana diagonal kang nyata.
\end{explain}

\begin{thm}
\ollabel{thm:nonenum-pownat}
$\Pow{\PosInt}$ ora !!{enumerable}.
\end{thm}

\begin{proof}
Kita nggunakake cara kang padha, yaiku nuduhake yen kanggo saben dhaptar
subhimpunan saka~$\PosInt$, ana subhimpunan $\PosInt$ kang ora bisa ana
ing dhaptar. Upamakna ana dhaptar subhimpunan saka~$\PosInt$:
\[
Z_{1}, Z_{2}, Z_{3}, \dots
\]
Saiki kita netepake himpunan $\overline{Z}$ saengga kanggo saben $n \in \PosInt$,
$n \in \overline{Z}$ yen lan mung yen $n \notin Z_{n}$:
\[
\overline{Z} = \Setabs{n \in \PosInt}{n \notin Z_n}
\]
Cetha yen $\overline{Z}$ iku himpunan wilangan wutuh positif, amarga
saben~$Z_n$ dianggep mangkono, mula $\overline{Z} \in
\Pow{\PosInt}$. Nanging $\overline{Z}$ ora bisa ana ing dhaptar.
Kanggo nuduhake iki, kita bakal mbuktekake yen kanggo saben $k \in \PosInt$,
$\overline{Z} \neq Z_k$.

Mula upamakna $k \in \PosInt$ sembarang. Kita wis netepake $\overline{Z}$
saengga kanggo saben $n \in \PosInt$, $n \in \overline{Z}$ yen lan mung yen
$n \notin Z_n$. Mligine, kanthi njupuk $n=k$, $k \in \overline{Z}$
yen lan mung yen $k \notin Z_k$. Nanging iki nuduhake
yen $\overline{Z} \neq Z_k$, amarga $k$ iku !!a{element} saka siji
himpunan nanging dudu anggota himpunan liyane. Dadi $\overline{Z}$ lan $Z_k$
duwe !!{element}s kang beda. Amarga $k$ sembarang, $\overline{Z}$
ora ana ing dhaptar $Z_1$, $Z_2$, \dots
\end{proof}

\begin{explain}
Bukti sadurunge ora nyebut diagonal, nanging kita bisa nganggep
ana diagonal yen nggambarake kaya mangkene: bayangna himpunan
$Z_1$, $Z_2$, \dots, ditulis ing larikan, kanthi saben
!!{element}~$j \in Z_i$ ditulis ing kolom kaping~$j$. Upamakna
papat himpunan kapisan ing dhaptar yaiku $\{1,2,3,\dots\}$, $\{2, 4, 6, \dots\}$,
$\{1,2,5\}$, lan $\{3,4,5,\dots\}$. Dadi wiwitan larikan:
\[
\begin{array}{r@{}rrrrrrr}
  Z_1 = \{ & \mathbf{1}, & 2, & 3, & 4, & 5, & 6, & \dots\}\\
  Z_2 = \{ &  & \mathbf{2}, &  & 4, &  & 6, & \dots\}\\
  Z_3 = \{ & 1, & 2, &  &  & 5\phantom{,} &  & \}\\
  Z_4 = \{ &  &  & 3, & \mathbf{4}, & 5, & 6, & \dots\}\\
  \vdots & & & & & \ddots
\end{array}
\]
Banjur $\overline{Z}$ iku himpunan kang diolehake kanthi ngetutake diagonal
mudhun, ora nglebokake wilangan kang katon ing diagonal, lan nglebokake
$j$ yen ana sela ing baris lan kolom kaping $j$.
Ing tuladha ndhuwur, kita ora nglebokake $1$ lan $2$, nglebokake~$3$,
ora nglebokake~$4$, lan sateruse.
\end{explain}

\begin{prob}
Tuduhna yen $\Pow{\Nat}$ iku !!{nonenumerable} nganggo argumen diagonal.
\end{prob}

\begin{prob}\label{sfr:siz:nen:prob:f-posint}
Tuduhna yen himpunan fungsi $f \colon \PosInt \to \PosInt$ iku
!!{nonenumerable} nganggo argumen diagonal kang eksplisit. Tegese,
tuduhna yen $f_1$, $f_2$, \dots, iku dhaptar fungsi lan saben $f_i\colon
\PosInt \to \PosInt$, mula ana $\overline{f}\colon \PosInt \to
\PosInt$ kang ora ana ing dhaptar iki.
\end{prob}

\end{document}
